\hypertarget{appendix-d49161831368}{}
\section*{Bitwise XOR}
\section{Bitwise XOR}\label{bitwise-xor}

\textbf{Curriculum depth:} Introduced · \textbf{First seen:} C001 · \textbf{Latest development:} C001

\subsection{Definition and mechanism}\label{definition-and-mechanism}

Toggles selected lanes and marks differences; applying the same mask twice restores the input.

XOR returns 1 when inputs differ. \texttt{0xAD\ \^{}\ 0x08\ =\ 0xA5} toggles visible; applying \texttt{0x08} again restores \texttt{0xAD}. XOR between old and new words marks changed positions. Reversibility alone does not provide secure encryption.

\subsection{Prerequisites}\label{prerequisites}

\begin{itemize}
\tightlist
\item
  \hyperref[appendix-24475e1596a8]{Boolean algebra} --- The logical domain is false/true.
\end{itemize}

\subsection{Used by}\label{used-by}

\begin{itemize}
\item
  \hyperref[appendix-ed338677de0a]{Checksums CRCs and hashes} --- uses this mechanism to define its representation, error or validation contract.
\item
  \hyperref[appendix-1652d27a8e90]{Half adder} --- uses bitwise xor as part of its representation or reasoning.
\end{itemize}

\subsection{Recall and next depth}\label{recall-and-next-depth}

Explain bitwise xor without the label, work the example above, then state the assumption whose removal would change the result.

\textbf{Pending:} parity, checksums, coding theory, and cryptographic constructions.

\subsection{Source and navigation}\label{source-and-navigation}

\href{Continuous\%20Notes/Computing\%20Foundations\%20-\%20From\%20Bits\%20to\%20C++\%20and\%20Systems\%20APIs\#From\%20a\%20signal\%20to\%20an\%20interpretation}{C001 teaching chapter} · \hyperref[appendix-fd090d982a15]{Part definitions and mechanisms} · \href{}{Dependency Map (02-Foundations-II, appendix-4ad47bdd269a)}
